\uuid{QhNq}
\titre{Matrice/Inversibilité}
\auteur{hanani ; mzari}
\organisation{exo7}
\licence{CC BY-NC-SA 4.0}
\datecreate{2020-01-20}
\niveau{L1}
\module{Algèbre}
\chapitre{Matrice}
\sousChapitre{Inverse, méthode de Gauss}
\difficulte{3}
\theme{Inverse d'une matrice, Matrice, Inversibilité}
\exo7qcmid{lille-2.matrices.027}

\contenu{
\begin{question}
\qclassification{108.05}
\qtags{motcle=Matrice/Inversibilité}
On considère la matrice : $A = \left(\begin{array}{rcc}
1&-1&-2\\0&1&1\\ -1&1&2\\ \end{array}\right)$. Quelles sont les assertions vraies ?
\begin{answers}
\bad{$A$ est inversible}.
\bad{$A^2$ est inversible}.
\bad{$A^3+A^2$ est inversible}.
\good{$A+\,^tA$ est inversible, où $^tA$ est la transposée de $A$}.
\end{answers}
\begin{explanations} $A$ n'est pas inversible, puisque les trois vesteurs colonnes sont linéairement dépendants.
\vskip0mm
Si $A^2$ est inversible, alors il existe une matrice  $B$ telle que $A^2B=I$, donc $A(AB)=I$ et donc $A$ est inversible,
ce qui est absurde.
\vskip0mm
Si $A^3+A^2$ est inversible, alors il existe une matrice $B$ telle que $(A^3+A^2)B=I$. On en déduit que $A[(A^2+A)B]=I$ et donc $A$ est inversible, ce qui est absurde.
\vskip0mm
$A+\, ^tA$ est inversible, puisque les vecteurs colonnes de cette matrice sont linéairement indépendants.
\end{explanations}
\end{question}
}
